\documentclass[10pt]{article}
\usepackage[T1]{fontenc}
\usepackage[utf8]{inputenc}
\usepackage[french]{babel}
\usepackage{amsmath,amssymb,mathrsfs,mathpazo}
\usepackage{geometry}
\geometry{a4paper,margin=2cm}
\usepackage{graphicx,array,tabularx,tabularray,colortbl,multirow,multicol,booktabs,enumitem,pifont,ulem,fancybox,diagbox,cancel}
\usepackage{pstricks,pstricks-add,pst-plot,pst-tree,pst-node,pst-text,pst-func,pst-math,pst-bspline,pst-3dplot,pst-eucl,pst-solides3d}
\let\pscalculate\undefined
\usepackage{pst-bezier}
\usepackage{tikz,pgfplots,tkz-tab}
\usetikzlibrary{arrows,arrows.meta,patterns,calc,positioning,trees,decorations.markings,shapes.geometric,shapes.misc,shapes.arrows,intersections,decorations.pathmorphing,decorations.pathreplacing,backgrounds,shadows,babel}
\pgfplotsset{compat=1.18}
\usepgfplotslibrary{fillbetween}
\usepackage[np]{numprint}
\usepackage{listings,makecell,pas-tableur,twoopt}
\let\barre\undefined
\usepackage[pstricks]{tablvar}
\definecolor{darkgreen}{RGB}{0,100,0}
\providecommand{\red}{\color{red}}
\providecommand{\blue}{\color{blue}}
\providecommand{\green}{\color{green}}
\providecommand{\black}{\color{black}}
\providecommand{\white}{\color{white}}
\providecommand{\gray}{\color{gray}}
\providecommand{\cyan}{\color{cyan}}
\providecommand{\magenta}{\color{magenta}}
\providecommand{\yellow}{\color{yellow}}
\providecommand{\violet}{\color{violet}}
\providecommand{\orange}{\color{orange}}
\providecommand{\euro}{\texteuro}
\providecommand{\EUR}{\texteuro}
\providecommand{\ieme}{\textsuperscript{e}}
\providecommand{\degres}{\ensuremath{^\circ}}
\providecommand{\degree}{\ensuremath{^\circ}}
\providecommand{\vv}[1]{\overrightarrow{#1}}
\providecommand{\Coord}[1]{#1}
\providecommand{\decosix}{}
\providecommand{\starredbullet}{\ensuremath{\bullet}}
\pagestyle{empty}
\setlength{\parindent}{0pt}
\usepackage{adjustbox}\DeclareUnicodeCharacter{2236}{:}\DeclareUnicodeCharacter{2190}{\ensuremath{\leftarrow}}
\providecommand{\R}{\mathbb{R}}
\providecommand{\N}{\mathbb{N}}
\providecommand{\Z}{\mathbb{Z}}
\providecommand{\Q}{\mathbb{Q}}
\providecommand{\D}{\mathbb{D}}
\providecommand{\C}{\mathbb{C}}
\providecommand{\e}{{\mathrm{e}}}
\providecommand{\vect}[1]{\overrightarrow{#1}}
\providecommand{\vectt}[1]{\overrightarrow{\mathrm{#1}}}
\providecommand{\vvt}[1]{\overrightarrow{\mathrm{#1}}}
\providecommand{\np}[2][]{{#2~#1}}
\providecommand{\vv}[1]{{\overrightarrow{#1}}}
\renewcommand{\d}{{\mathrm{d}}}
\providecommand{\degres}{{^\circ}}
\providecommand{\degre}{{^\circ}}
\providecommand{\blue}{\color{blue}}
\providecommand{\red}{\color{red}}
\providecommand{\enskip}{\quad}
\providecommand{\llbracket}{\lbrack\!\lbrack}
\providecommand{\rrbracket}{\rbrack\!\rbrack}
\providecommand{\euro}{{\text{€}}}
\providecommand{\no}{{\text{n°}}}
\providecommand{\interval}[2]{{[#1;#2]}}
\providecommand{\intervFO}[2]{{[#1;#2[}}
\providecommand{\dotfill}{\ldots}

\providecommand{\ds}{\displaystyle}
\providecommand{\textding}[1]{#1}
\providecommand{\Oij}{\(\mathrm{O};\vec{\imath},\vec{\jmath}\)}
\providecommand{\Oijk}{\(\mathrm{O};\vec{\imath},\vec{\jmath},\vec{k}\)}
\providecommand{\Ouv}{\(\mathrm{O};\vec{u},\vec{v}\)}
\providecommand{\intervFF}[2]{\ensuremath{[#1;#2]}}
\renewcommand{\euro}{\ifmmode\text{\texteuro}\else\texteuro\fi}
\DeclareUnicodeCharacter{00B0}{\ifmmode^{\circ}\else\textdegree\fi}
\DeclareUnicodeCharacter{20AC}{\ifmmode\text{\texteuro}\else\texteuro\fi}
% Énoncés : Éducation nationale. Transcription et corrigés : APMEP, Denis Vergès et contributeurs. Découpage et classement : CoopMaths, Sébastien Lozano, Éric Elter et contributeurs. Adaptation de la présentation : Maths Explorer.
% CC BY-SA : https://coopmaths.fr/www/about/
% https://coopmaths.fr/alea/static/bac/2025/tex/bac_2025_06_sujet1_polynesie_4_cor.tex
\begin{document}
\noindent\textit{APMEP / CoopMaths. CC BY-SA. Présentation adaptée par Maths Explorer.}\par\medskip
\medskip

\begin{enumerate}
	\item \textbf{Affirmation 1 : Vraie.}

	
	
	

	On reconnaît une forme présente dans le cours : les équations différentielles de la forme $y' = ay + b$, avec $a$ et $b$ réels non nuls ont pour solution les fonctions de la forme : \quad $x \longmapsto k\e^{ax} - \dfrac{b}{a}$,\quad où $k$ est un réel quelconque.

	Ici, on applique cette propriété avec $a = \dfrac{1}{2}$ et $b = 4$, ce qui donne bien l'équation différentielle $y' = \dfrac{1}{2}y + 4$, ayant pour solution les fonction de la forme :\quad$x \longmapsto k\e^{\frac{1}{2}x} - \dfrac{4}{\frac{1}{2}}$ soit $x \longmapsto k\e^{\frac{1}{2}x} - 8$, avec $k\in \R$.

	
	
		
		
		
		
		
		
		
			
			
			
			
			
			
		
		
		
		

	\item \textbf{Affirmation 2 : Vraie.}

	Une équipe de volley-ball telle que décrite est constituée d'un couple de deux éléments. Le premier élément de ce couple est un ensemble de 3 élèves filles distinctes (donc pas de répétition), qui n'ont pas de rôle particulier, donc on parle d'un ensemble, et pas d'une liste ou d'un triplet (donc pas de notion d'ordre) : on va utiliser des combinaisons.

	3 filles choisies parmi les 18 de la classe : il y a $\displaystyle \binom{18}{3} = 816$ façons de choisir les trois filles qui constitueront ce premier ensemble.

	Le deuxième élément du couple sera un ensemble de trois élèves garçons, choisis avec les même principes.

	3 garçons choisis parmi les 14 de la classe : il y a $\displaystyle \binom{14}{3} = 364$ façons de choisir les trois garçons qui constitueront ce second ensemble.

	Finalement, par principe multiplicatif, une équipe est un ensemble de trois filles choisi parmi les 816 ensembles possibles, associé à un ensemble de garçons parmi les 364 ensembles possibles : il y a donc $\quad816 \times 364 = \np{297024}$ équipes différentes possibles.

	\item \textbf{Affirmation 3 : Vraie.}

	Pour tout $n$ entier naturel, on a :

	$\begin{aligned}[t]
		-1 \leqslant \cos(n) \leqslant 1 &\implies 1\leqslant 2 + \cos(n) \leqslant 3\\
		&\implies \dfrac{1}{3} \leqslant \dfrac{1}{2+\cos(n)} \leqslant \dfrac{1}{1}\quad\text{fonction inverse est décroissante sur } \R_{+}^{*}\\
		&\implies \dfrac{n}{3} \leqslant \dfrac{n}{2 + \cos(n)}\quad \text{car } n\geqslant 0\\
		&\implies \dfrac{n}{3} \leqslant v_n
	\end{aligned}$

	Par limite du quotient, on a immédiatement : \quad$\lim\limits_{n \to +\infty} \dfrac{n}{3} = +\infty$,

	donc, par comparaison, \quad $\lim\limits_{n \to +\infty} v_n = +\infty$,\quad ce qui est l'affirmation 3.

	\item \textbf{Affirmation 4 : Fausse.}

	On a : $\vect{\mathrm{AB}}\begin{pmatrix}4\\-2\\6\end{pmatrix}$ et $\vect{\mathrm{AC}}\begin{pmatrix}1\\0\\1\end{pmatrix}$,
	donc, comme le repère est orthonormé, on a :

	$\vect{\mathrm{AB}}\cdot\vect{\mathrm{AC}} = 4\times 1 + (-2)\times 0 + 6 \times 1 = 10$,\quad la première partie de l'affirmation est donc vraie.

	Cependant, on peut aussi calculer les normes et en déduire le cosinus de l'angle $\widehat{\mathrm{BAC}}$, défini par les deux vecteurs.

	$\mathrm{AB} = \sqrt{4^2 + (-2)^2 + 6^2} = \sqrt{16 + 4 + 36} = \sqrt{56} = \sqrt{4 \times 14} = 2\sqrt{14}$.

	$\mathrm{AC} = \sqrt{1^2 + 0^2 + 1^2} = \sqrt{1 + 0 + 1} = \sqrt{2}$.

	On a donc, avec la définition trigonométrique du produit scalaire :

	$ \vect{\mathrm{AB}}\cdot\vect{\mathrm{AC}} = \mathrm{AB}\times \mathrm{AC}\times \cos\left(\widehat{\mathrm{BAC}}\right) = 2\sqrt{14}\times \sqrt{2} \times \cos\left(\widehat{\mathrm{BAC}}\right) = 4\sqrt{7}\cos\left(\widehat{\mathrm{BAC}}\right)$

	En égalant ces deux expressions du produit scalaire, on déduit :

	$10 = 4\sqrt{7}\cos\left(\widehat{\mathrm{BAC}}\right)$\quad qui équivaut à :\quad $\cos\left(\widehat{\mathrm{BAC}}\right) = \dfrac{10}{4\sqrt{7}} = \dfrac{5\sqrt{7}}{14}$.

	Il se trouve que le cosinus de 30\ensuremath{^\circ}{} est connu, il vaut $\dfrac{\sqrt{3}}{2}$ et est donc différent de celui trouvé ici : l'angle ne mesure dont pas 30\ensuremath{^\circ}.

\textbf{Remarque :} une valeur approchée d'une mesure de l'angle peut être obtenue à la calculatrice :

$\widehat{\mathrm{BAC}} = \arccos\left(\dfrac{5\sqrt{7}}{14}\right)\approx 19$\ensuremath{^\circ}.

\item \textbf{Affirmation 5 : Vraie.}

En effet, pour tout x réel strictement positif, on a :\quad $ h''(x) = x\ln(x) - 3x = x\big(\ln(x) - 3\big)$

$x$ étant strictement positif, le signe de $h''(x)$ est donc le signe de $\ln(x)-3$.

	$\begin{aligned}
		x \in \big[\e^{3}~;~+\infty\big[ &\implies \e^{3} \leqslant x\\
		&\implies \ln\big(\e^{3}\big) \leqslant \ln(x)\quad\text{car ln est croissante sur }\R*+\\
		&\implies 3\leqslant \ln(x)\\
		&\implies 0 \leqslant \ln(x) - 3\\
		&\implies 0\leqslant h''(x)
	\end{aligned}$

	Sur l'intervalle $\big[\e^{3}~;~+\infty\big[$, la dérivée seconde de $h$ est donc à valeurs positives, donc la dérivée première de $h$ est croissante, et donc $h$ elle-même est bien convexe sur cet intervalle.
\end{enumerate}

\end{document}
